% Extensión cuántica de la teoría general de la medición.
% Se integra en el mismo capítulo fuente de 01_medida_ambivalencia.

\section{Condicionamiento de secciones y actualización cuántica}
\label{sec:medicion-cuantica-genealogica-rev2}
\index{medición!condicionamiento de secciones}
\index{instrumento cuántico!realización genealógica}

La teoría general precedente sitúa la medición en el estado conjunto de
sistema, lector y aparato. La realización cuántica añade una correspondencia
precisa entre tres niveles: la restricción del espacio de secciones, la
proyección espectral y el instrumento completamente positivo. La evolución
conjunta pertenece al acoplamiento; la publicación del resultado pertenece al
lector; la actualización pertenece al condicionamiento por la fibra leída.

\subsection{Contexto de medición \(\mathfrak M_a=(U_a,q_a,\mathcal C_a,m_0)\), aplicación de salida y fibra condicionada \(\Omega_{a,y}\)}

Sean \((\Omega_S,\Sigma_S)\) y \((\Omega_M,\Sigma_M)\) los espacios
medibles de estados genealógicos del sistema y del aparato. Un contexto
\(a\) se especifica mediante
\begin{equation}
 \mathfrak M_a=(U_a,q_a,\mathcal C_a,m_0),
 \label{eq:contexto-medida-rev2}
\end{equation}
donde \(m_0\in\Omega_M\) es la preparación del aparato,
\[
 U_a:\Omega_S\times\Omega_M
 \longrightarrow\Omega_S\times\Omega_M
\]
es un acoplamiento medible, \((Y_a,\Sigma_{Y_a})\) es el espacio medible de
resultados, \(q_a:\Omega_M\to Y_a\) es un lector medible del puntero y
\(\mathcal C_a\) declara la compatibilidad y el horizonte de prolongación.
Cuando se invoque un régimen reversible se exigirá, además, que \(U_a\) sea
biyectivo y bimedible; en la realización de Hilbert, la condición
correspondiente es la unitariedad de \(\widehat U_a\).
La salida es
\begin{equation}
 \operatorname{Out}_a(x)
 =q_a\!\left(\operatorname{pr}_M U_a(x,m_0)\right).
 \label{eq:salida-contexto-medida-rev2}
\end{equation}

Sea \(\Omega_a\in\Sigma_S\) el conjunto de secciones compatibles con
la preparación y el contexto. La aplicación \(\operatorname{Out}_a\) es
medible por composición. Para \(y\in Y_a\) con
\(\{y\}\in\Sigma_{Y_a}\), se define la fibra medible
\begin{equation}
 \Omega_{a,y}
 :=\{x\in\Omega_a:\operatorname{Out}_a(x)=y\}.
 \label{eq:fibra-resultado-medida-rev2}
\end{equation}

\begin{definition}[Actualización genealógica]
La actualización determinada por el resultado \(y\) es la restricción
\[
 \Omega_a\longmapsto\Omega_{a,y}.
\]
En profundidad finita conserva exactamente los prefijos que admiten alguna
extensión en \(\Omega_{a,y}\).
\end{definition}

Para un prefijo \(s\in S_n\), los futuros compatibles posteriores a la
lectura forman
\[
 \operatorname{Fut}_{a,y}(s)
 =\{x\in\Omega_{a,y}:x|_n=s\}.
\]
La actualización modifica la condición de prolongación y conserva el prefijo
ya recorrido junto con el registro del acoplamiento.

\begin{proposition}[Probabilidad condicionada sobre la fibra]
\label{prop:probabilidad-condicionada-medida-rev2}
Sea \(\mu_a\) una probabilidad sobre el espacio medible restringido
\((\Omega_a,\Sigma_S|_{\Omega_a})\), sea
\(B\in\Sigma_S|_{\Omega_a}\) y supóngase
\(\mu_a(\Omega_{a,y})>0\). Entonces
\begin{equation}
 \mu_{a,y}(B)
 =\frac{\mu_a(B\cap\Omega_{a,y})}
        {\mu_a(\Omega_{a,y})}
 \label{eq:bayes-genealogico-rev2}
\end{equation}
define una probabilidad concentrada en la fibra del resultado:
\(\mu_{a,y}(\Omega_{a,y})=1\).
\end{proposition}

\begin{proof}
La aplicación es no negativa y numerablemente aditiva por las propiedades de
\(\mu_a\). La normalización se obtiene sustituyendo
\(B=\Omega_{a,y}\) en \eqref{eq:bayes-genealogico-rev2}.
\end{proof}

\subsection{Realización espectral del aparato}

Sea \(\{|r\rangle\}\) una base ortonormal adaptada al observable del sistema,
sea \(|m_0\rangle\) la preparación del aparato y sea
\(\{|m_r\rangle\}\) una familia ortonormal de estados de puntero. La
representación unitaria del acoplamiento satisface
\begin{equation}
 \widehat U_a(|r\rangle\otimes|m_0\rangle)
 =|r\rangle\otimes|m_r\rangle.
 \label{eq:premedicion-espectral-rev2}
\end{equation}
Para \(|\psi\rangle=\sum_r c_r|r\rangle\), la linealidad produce
\[
 \widehat U_a(|\psi\rangle\otimes|m_0\rangle)
 =\sum_r c_r|r\rangle\otimes|m_r\rangle.
\]
El acoplamiento construye así una correlación entre sistema y puntero. La
lectura selectiva del resultado \(r\) se representa, para un estado densidad
con \(\operatorname{tr}(P_r\rho)>0\), por la regla de
Lüders~\cite{Lueders1950},
\begin{equation}
 \rho\longmapsto
 \rho_r=\frac{P_r\rho P_r}{\operatorname{tr}(P_r\rho)}.
 \label{eq:luders-genealogico-rev2}
\end{equation}

\begin{theorem}[Correspondencia entre fibras y proyecciones]
\label{thm:fibra-proyeccion-medida-rev2}
Supóngase que cada fibra \(\Omega_{a,y}\) es una unión finita de cilindros
disjuntos de una profundidad común \(n\), que la realización
\(\mathcal Q\) lleva esos cilindros a subespacios ortogonales y que \(P_y\)
es la suma de sus proyectores cilíndricos. Sea
\(\mu_a(\Omega_{a,y})>0\) y sea \(\Psi_n(\mu_a)\) el estado construido por
\eqref{eq:estado-amplitudes-medida-rev2} a partir de esa misma medida. Entonces
\[
 \mu_a(\Omega_{a,y})=\|P_y\Psi_n(\mu_a)\|^2,
\]
y la restricción a \(\Omega_{a,y}\) se representa por la proyección
normalizada
\(P_y\Psi_n(\mu_a)/\|P_y\Psi_n(\mu_a)\|\). En la carta de matrices
densidad, si
\(\operatorname{tr}(\rho P_y)>0\), sea \(\mathcal A_n\) el álgebra de
sucesos cilíndricos de profundidad \(n\) y supóngase que \(B\mapsto P_B\) y
\(P_y\) pertenecen al mismo PVM y que
\begin{equation}
 \mu_a(B)=\operatorname{tr}(\rho P_B)
 \qquad\text{para todo }B\in\mathcal A_n.
 \label{eq:equivalencia-medida-pvm-rev2}
\end{equation}
Entonces la misma probabilidad condicionada se representa por
\eqref{eq:luders-genealogico-rev2}: para todo \(B\in\mathcal A_n\),
\[
 \mu_{a,y}(B)=\operatorname{tr}(\rho_yP_B).
\]
\end{theorem}

\begin{proof}
La fibra es una unión disjunta de cilindros y \(P_y\) es la suma ortogonal de
sus proyectores. El \cref{thm:born-cilindrico-rev2} identifica la norma
cuadrática de cada componente con su medida; la aditividad da la igualdad.
Tras la proyección, los coeficientes de cada cilindro se dividen por
\(\sqrt{\mu_a(\Omega_{a,y})}\); sus cuadrados son, por tanto, las
probabilidades condicionadas de \eqref{eq:bayes-genealogico-rev2}.
En la carta de densidades, la pertenencia al mismo PVM da
\(P_BP_y=P_{B\cap\Omega_{a,y}}\). Por ciclicidad de la traza y por
\eqref{eq:equivalencia-medida-pvm-rev2},
\[
 \operatorname{tr}(\rho_yP_B)
 =\frac{\operatorname{tr}(\rho P_{B\cap\Omega_{a,y}})}
        {\operatorname{tr}(\rho P_y)}
 =\frac{\mu_a(B\cap\Omega_{a,y})}{\mu_a(\Omega_{a,y})}.
\]
\end{proof}

\subsection{Instrumentos inducidos por el acoplamiento}

Para una base de punteros refinada
\(\{|m_{y,\alpha}\rangle\}_{y,\alpha}\), los operadores
\begin{equation}
 M_{y,\alpha}=\langle m_{y,\alpha}|\widehat U_a|m_0\rangle,
 \qquad
 \sum_{y,\alpha} M_{y,\alpha}^\dagger M_{y,\alpha}=I
 \label{eq:kraus-acoplamiento-genealogico-rev2}
\end{equation}
forman una representación de Kraus~\cite{Kraus1983}. El efecto
\(E_y=\sum_\alpha M_{y,\alpha}^\dagger M_{y,\alpha}\) determina la
probabilidad y, cuando \(p(y\mid a,\rho)>0\), el estado posterior:
\begin{equation}
 p(y\mid a,\rho)=\operatorname{tr}(\rho E_y),
 \qquad
 \rho_{a,y}
 =\frac{\sum_\alpha M_{y,\alpha}\rho M_{y,\alpha}^\dagger}
 {p(y\mid a,\rho)}.
 \label{eq:instrumento-kraus-genealogico-rev2}
\end{equation}
La transformación no selectiva es
\[
 \mathcal E_a(\rho)=
 \sum_{y,\alpha}M_{y,\alpha}\rho M_{y,\alpha}^\dagger.
\]

\begin{proposition}[Acoplamiento unitario y canal reducido]
\label{prop:acoplamiento-canal-reducido-rev2}
La dinámica conjunta \(\widehat U_a\) y la preparación \(|m_0\rangle\)
inducen un canal completamente positivo y preservador de la traza sobre el
sistema. La descomposición por punteros induce el instrumento
\(\mathcal I_{a,y}(\rho)=
\sum_\alpha M_{y,\alpha}\rho M_{y,\alpha}^\dagger\). Recíprocamente, todo canal
completamente positivo admite una dilatación por un espacio auxiliar y una
isometría.
\end{proposition}

\begin{proof}
La identidad de completitud en
\eqref{eq:kraus-acoplamiento-genealogico-rev2} preserva la traza, y la forma
de Kraus demuestra positividad completa. La afirmación recíproca es el
teorema de dilatación de Stinespring~\cite{Stinespring1955}.
\end{proof}

El espacio auxiliar de la dilatación puede realizar linealmente la memoria y
el registro transportados por TPK una vez declarado un encaje que entrelace
esa memoria con el espacio auxiliar. Stinespring garantiza la dilatación del
canal; no identifica por sí solo ambos objetos. La regla de Lüders es el caso
proyectivo fino \(M_{y,1}=P_y\); los lectores degenerados y los aparatos no
ideales se describen mediante familias \(M_{y,\alpha}\).

\subsection{Observador interno y completitud de la lectura}

Un observador realizado es un subsistema capaz de preparar un contexto,
acoplarse a una frontera, conservar un registro distinguible y utilizarlo
como condición de prolongación. La secuencia
\begin{equation}
 (x,m_0)
 \xrightarrow{\ U_a\ }(x',m')
 \xrightarrow{\ q_a\ }y
 \xrightarrow{\ \mathrm{cond}\ }\Omega_{a,y}
 \label{eq:secuencia-observador-rev2}
\end{equation}
sitúa el marco de lectura y la memoria del observador dentro del estado
ampliado.

El principio de Ambivalencia completa esta descripción. Las lecturas areal y
volumétrica producen cocientes \(q_A,q_V\), mientras la memoria \(T\)
conserva el transporte entre ambas. Cuando
\[
 \Xi:x\longmapsto(q_A(x),q_V(x),T(x))
\]
es inyectiva módulo el grupo de calibre declarado, la pareja de observables y
la memoria reconstruyen la clase del estado completo.

\begin{theorem}[Descomposición genealógica de la medición]
\label{thm:descomposicion-medicion-rev2}
Bajo las hipótesis del
\cref{thm:fibra-proyeccion-medida-rev2,prop:acoplamiento-canal-reducido-rev2}, la
transición entre evolución y resultado se descompone, en el régimen en que
\(U_a\) es biyectivo y bimedible o \(\widehat U_a\) es unitario, en:
\[
 \text{acoplamiento reversible}
 \longrightarrow
 \text{lectura del puntero}
 \longrightarrow
 \text{condicionamiento de secciones}.
\]
Esta descomposición posee dos regímenes:
\begin{enumerate}
\item En una medición cilíndrica proyectiva, \(M_{y,1}=P_y\) y coinciden la
restricción de fibra, la proyección normalizada y la actualización de Lüders:
\[
 \Omega_a\longmapsto\Omega_{a,y},
 \qquad
 \Psi_n\longmapsto\frac{P_y\Psi_n}{\|P_y\Psi_n\|},
 \qquad
 \rho\longmapsto
 \frac{P_y\rho P_y}{\operatorname{tr}(\rho P_y)}.
\]
\item Para un instrumento general, el estado condicionado es
\[
 \rho\longmapsto
 \frac{\sum_\alpha M_{y,\alpha}\rho M_{y,\alpha}^\dagger}
 {\operatorname{tr}(\rho E_y)}.
\]
En una dilatación, esta actualización procede de una proyección del puntero
sobre el sistema compuesto con el aparato. Su correspondencia con
\(\Omega_{a,y}\) requiere un morfismo declarado que entrelace la fibra del
lector, el efecto \(E_y\) y el instrumento \(\mathcal I_{a,y}\).
\end{enumerate}
\end{theorem}

\begin{proof}
La primera flecha realiza el acoplamiento; la tercera restringe el espacio de
secciones a la fibra definida en \eqref{eq:fibra-resultado-medida-rev2}. El
\cref{thm:fibra-proyeccion-medida-rev2} demuestra la coincidencia de las tres
formas en el régimen proyectivo. En el régimen general,
\eqref{eq:instrumento-kraus-genealogico-rev2} demuestra la actualización del
instrumento y Stinespring proporciona su realización dilatada. El morfismo de
entrelazamiento adicional es precisamente el dato que transporta esa
actualización de vuelta al espacio de secciones.
\end{proof}

\begin{statusbox}[title={Resultado y residencia}]
En el modelo declarado, la actualización cuántica es la realización
espectral de la restricción de futuros compatibles. La reversibilidad se
decide por la bimedibilidad de \(U_a\) o la unitariedad de
\(\widehat U_a\); la pérdida informativa, en las fibras del lector o en un
borrado posterior; y el estado condicionado, en \(\Omega_{a,y}\). El coste
termodinámico de reinicialización y el Landauer nulo de la evolución
reversible permanecen en el capítulo de información y entropía.
\end{statusbox}
